\section{Task Definition}

Given a medical question $q$, a biomedical corpus $C$, and optional knowledge sources $K$, the system returns a ranked evidence list $E$ and optionally an answer $a$. The primary task is evidence retrieval and reranking; answer generation is secondary.

For each question, a first-stage retriever returns a candidate pool $D_q = \{d_1,\ldots,d_M\}$ from $C$. A reranker then estimates an improved ordering over $D_q$ using retrieval scores, biomedical semantic scores, concept and entity features, and local hypergraph features. Evaluation emphasizes Recall@$k$, MRR@$k$, nDCG@$k$, and evidence coverage.

The formal output is:
\[
  (a, E, P),
\]
where $a$ is optional in the current evidence-focused setting, $E$ is the ranked supporting evidence list, and $P$ is an optional interpretable entity, concept, or hypergraph path used for case analysis.

